\MainChapter{Justificación}{\topBanner}\label{cap:justificacion}
\vspace*{0.7cm}

En el trabajo de \citet{Felts2023} se evidencia que, en los últimos años, las instituciones educativas han aumentado el uso de recursos digitales y la colaboración entre usuarios, lo que evidencia la necesidad de contar con mecanismos centralizados para gestionar de manera eficiente a los usuarios y sus accesos. Cuando cada servicio administra sus propias cuentas y permisos, se generan mayores dificultades operativas, aumentan los costos de gestión y aparecen riesgos de seguridad por la falta de políticas unificadas de autenticación y autorización \citep{Pohn2020}. En este sentido, la autenticación federada permite reducir de manera significativa la carga administrativa de la gestión de cuentas, a la vez que incrementa la seguridad general del sistema y mejora la experiencia del usuario al integrar su cuenta como punto central de acceso \citep{Prout2019}. En este contexto el grupo de investigación \ac{GRID} tiene en los servicios de directorio una oportunidad para dar estructura a los procesos de gestión de cuentas de usuarios y permisos sobre los recursos hardware y software disponibles.
